\section{Hasil Peramalan}
Berdasarkan hasil peramalan yang telah dilakukan dengan ketiga pembobotan itu diperoleh yang paling baik yaitu Bobot Seragam dengan model STARIMA $(0,0,0) x (0,1,1_{0})_{12}$.
Hasil yang diperoleh, dapat disimpulkan bahwa model STARIMA mengungguli model SARIMA univariat. Keunggulan model STARIMA dibandingkan model SARIMA univariat dapat disebabkan oleh penyertaan informasi    spasial, yaitu efek ketetanggaan, dalam bentuk matriks bobot spasial.
\begin{table}[H]
    \centering
    \scriptsize
    \caption{Perbandingan RMSE SARIMA dan STARIMA per Wilayah}
    \label{table:perbandingan-rmse-arima-starima}
    \renewcommand{\arraystretch}{1.25}

    \begin{tabular}{l l c}
        \toprule
        \textbf{Wilayah} & \textbf{Model} & \textbf{RMSE} \\
        \midrule

        \multirow{2}{*}{Barat}
            & SARIMA$(1,0,0) \times (0,1,0)_{12}$    & 5.082 \\
            & STARIMA$(0,0,0) \times (0,1,1_{0})_{12}$        & 2.962 \\
        \midrule

        \multirow{2}{*}{Selatan}
            & SARIMA$(1,0,0) \times (0,1,0)_{12}$    & 4.680 \\
            & STARIMA$(0,0,0) \times (0,1,1_{0})_{12}$   & 3.921 \\
        \midrule

        \multirow{2}{*}{Tengah}
            & SARIMA$(1,0,0) \times (0,1,0)_{12}$    & 4.434 \\
            & STARIMA$(0,0,0) \times (0,1,1_0)_{12}$        & 4.092 \\
        \midrule

        \multirow{2}{*}{Timur}
            & SARIMA$(1,0,0) \times (0,1,0)_{12}$    & 4.643 \\
            & STARIMA$(0,0,0) \times (0,1,1_0)_{12}$        & 1.331 \\
        \midrule

        \multirow{2}{*}{Utara}
            & SARIMA$(1,0,0) \times (0,1,0)_{12}$    & 5.016 \\
            & STARIMA$(0,0,0) \times (0,1,1_0)_{12}$        & 2.992 \\
        \bottomrule
    \end{tabular}
\end{table}

Tabel~\ref{table:perbandingan-rmse-arima-starima} ini menunjukkan nilai RMSE dari model SARIMA univariat dan model STARIMA untuk masing - masing wilayah.
Berdasarkan hasil peramalan, model STARIMA secara konsisten memberikan nilai kesalahan (RMSE) yang lebih rendah dibandingkan SARIMA.
Hal ini mengindikasikan bahwa penyertaan informasi spasial melalui matriks bobot spasial mampu meningkatkan akurasi model secara signifikan.
Dengan demikian, STARIMA terbukti lebih unggul dalam memodelkan pola curah hujan antar wilayah yang saling berhubungan secara geografis.